% Part: incompleteness
% Chapter: representability-in-q
% Section: representing-relations

\documentclass[../../../include/open-logic-section]{subfiles}

\begin{document}

\olfileid{inc}{req}{rel}

\olsection{تمثيل العلاقات}

نحدّ الآن قابلية التمثيل في حق \emph{العلاقات}، بعد أن حدّدناها في حق الدوال.

\begin{defn}
\ollabel{defn:representing-relations} لتكن
$\OLId{latin-looped}{R}(\OLId{latin-ordinary}{x}_\OLMathNumeral{0},\dots,\OLId{latin-ordinary}{x}_\OLId{latin-ordinary}{k})$ علاقة على الأعداد الطبيعية.
نقول إنها {\em قابلة للتمثيل في $\Th{Q}$} إذا أمكن اختيار صيغة
$!A_\OLId{latin-looped}{R}(\OLId{latin-ordinary}{x}_\OLMathNumeral{0},\dots,\OLId{latin-ordinary}{x}_\OLId{latin-ordinary}{k})$ تستوفي الشرطين الآتيين:
إذا صدقت $\OLId{latin-looped}{R}(\OLId{latin-ordinary}{n}_\OLMathNumeral{0},\dots,\OLId{latin-ordinary}{n}_\OLId{latin-ordinary}{k})$ أثبتت $\Th{Q}$
$!A_\OLId{latin-looped}{R}(\num{\OLId{latin-ordinary}{n}_\OLMathNumeral{0}},\dots,\num{\OLId{latin-ordinary}{n}_\OLId{latin-ordinary}{k}})$،
وإذا كذبت $\OLId{latin-looped}{R}(\OLId{latin-ordinary}{n}_\OLMathNumeral{0},\dots,\OLId{latin-ordinary}{n}_\OLId{latin-ordinary}{k})$ أثبتت $\Th{Q}$
$\lnot !A_\OLId{latin-looped}{R}(\num{\OLId{latin-ordinary}{n}_\OLMathNumeral{0}},\dots,\num{\OLId{latin-ordinary}{n}_\OLId{latin-ordinary}{k}})$.
\end{defn}

\begin{thm}
\ollabel{thm:representing-rels} قابلية العلاقة للتمثيل في
$\Th{Q}$ تكافئ قابليتها للحساب.
\end{thm}

\begin{proof}
أما من التمثيل إلى الحساب، فلتكن الصيغة $!A_\OLId{latin-looped}{R}(\OLId{latin-ordinary}{x}_\OLMathNumeral{0},\dots,\OLId{latin-ordinary}{x}_\OLId{latin-ordinary}{k})$ ممثلةً
لـ$\OLId{latin-looped}{R}(\OLId{latin-ordinary}{x}_\OLMathNumeral{0},\dots,\OLId{latin-ordinary}{x}_\OLId{latin-ordinary}{k})$. نحسب $\OLId{latin-looped}{R}$ بالإجراء الآتي:
إذا أعطيت $\OLId{latin-ordinary}{n}_\OLMathNumeral{0}$، \dots،~$\OLId{latin-ordinary}{n}_\OLId{latin-ordinary}{k}$، أجرينا البحث عن برهان
$!A_\OLId{latin-looped}{R}(\num{\OLId{latin-ordinary}{n}_\OLMathNumeral{0}},\dots,\num{\OLId{latin-ordinary}{n}_\OLId{latin-ordinary}{k}})$ وعن برهان
$\lnot !A_\OLId{latin-looped}{R}(\num{\OLId{latin-ordinary}{n}_\OLMathNumeral{0}},\dots,\num{\OLId{latin-ordinary}{n}_\OLId{latin-ordinary}{k}})$ بالتوازي.
وبفرض التمثيل ينتهي أحد البحثين، ولا يجتمع البرهانان لاتساق النظرية.
فإن وُجد الأول كان الجواب «نعم»، وإن وُجد الآخر كان الجواب «لا».

وأما من الحساب إلى التمثيل، فنفرض أن $\OLId{latin-looped}{R}(\OLId{latin-ordinary}{x}_\OLMathNumeral{0},\dots,\OLId{latin-ordinary}{x}_\OLId{latin-ordinary}{k})$ قابلة للحساب.
وهذا، بحسب التعريف، هو قابلية الدالة $\Char{\OLId{latin-looped}{R}}(\OLId{latin-ordinary}{x}_\OLMathNumeral{0},\dots,\OLId{latin-ordinary}{x}_\OLId{latin-ordinary}{k})$ للحساب.
ومن \olref[int]{thm:representable-iff-comp} نجد لتمثيل $\Char{\OLId{latin-looped}{R}}$
صيغةً $!A_{\Char{\OLId{latin-looped}{R}}}(\OLId{latin-ordinary}{x}_\OLMathNumeral{0},\dots,\OLId{latin-ordinary}{x}_\OLId{latin-ordinary}{k},\OLId{latin-ordinary}{y})$.
فنجعل $!A_\OLId{latin-looped}{R}(\OLId{latin-ordinary}{x}_\OLMathNumeral{0},\dots,\OLId{latin-ordinary}{x}_\OLId{latin-ordinary}{k})$ هي الصيغة
$!A_{\Char{\OLId{latin-looped}{R}}}(\OLId{latin-ordinary}{x}_\OLMathNumeral{0},\dots,\OLId{latin-ordinary}{x}_\OLId{latin-ordinary}{k},\num{\OLMathNumeral{1}})$.
وخذ الآن $\OLId{latin-ordinary}{n}_\OLMathNumeral{0}$، \dots،~$\OLId{latin-ordinary}{n}_\OLId{latin-ordinary}{k}$ كيفما كانت.
إذا صدقت $\OLId{latin-looped}{R}(\OLId{latin-ordinary}{n}_\OLMathNumeral{0},\dots,\OLId{latin-ordinary}{n}_\OLId{latin-ordinary}{k})$، كانت
$\Char{\OLId{latin-looped}{R}}(\OLId{latin-ordinary}{n}_\OLMathNumeral{0},\dots,\OLId{latin-ordinary}{n}_\OLId{latin-ordinary}{k})=\OLMathNumeral{1}$؛ فتثبت $\Th{Q}$
$!A_{\Char{\OLId{latin-looped}{R}}}(\num{\OLId{latin-ordinary}{n}_\OLMathNumeral{0}},\dots,\num{\OLId{latin-ordinary}{n}_\OLId{latin-ordinary}{k}},\num{\OLMathNumeral{1}})$،
أي تثبت $\Th{Q}$ الصيغة $!A_\OLId{latin-looped}{R}(\num{\OLId{latin-ordinary}{n}_\OLMathNumeral{0}},\dots,\num{\OLId{latin-ordinary}{n}_\OLId{latin-ordinary}{k}})$.
وأما إذا كذبت $\OLId{latin-looped}{R}(\OLId{latin-ordinary}{n}_\OLMathNumeral{0},\dots,\OLId{latin-ordinary}{n}_\OLId{latin-ordinary}{k})$، كانت $\Char{\OLId{latin-looped}{R}}(\OLId{latin-ordinary}{n}_\OLMathNumeral{0},\dots,\OLId{latin-ordinary}{n}_\OLId{latin-ordinary}{k})=\OLMathNumeral{0}$،
فتثبت $\Th{Q}$:
\[
\lforall[\OLId{latin-ordinary}{y}][(!A_{\Char{\OLId{latin-looped}{R}}}(\num{\OLId{latin-ordinary}{n}_\OLMathNumeral{0}}, \dots, \num{\OLId{latin-ordinary}{n}_\OLId{latin-ordinary}{k}}, \OLId{latin-ordinary}{y}) \lif \OLId{latin-ordinary}{y} =
  \num{\OLMathNumeral{0}})].
\]
ثم إن $\Th{Q}$ تثبت $\eq/[\num{\OLMathNumeral{0}}][\num{\OLMathNumeral{1}}]$؛
فلا بد أن تثبت $\Th{Q}$
$\lnot !A_{\Char{\OLId{latin-looped}{R}}}(\num{\OLId{latin-ordinary}{n}_\OLMathNumeral{0}},\dots,\num{\OLId{latin-ordinary}{n}_\OLId{latin-ordinary}{k}},\num{\OLMathNumeral{1}})$،
وهذا هو $\lnot !A_\OLId{latin-looped}{R}(\num{\OLId{latin-ordinary}{n}_\OLMathNumeral{0}},\dots,\num{\OLId{latin-ordinary}{n}_\OLId{latin-ordinary}{k}})$ المطلوب.
\end{proof}

\begin{prob}
استخرج من قابلية $\OLId{latin-looped}{R}$ للتمثيل في~$\Th{Q}$
قابليةَ $\Char{\OLId{latin-looped}{R}}$ للتمثيل فيها.
\end{prob}

\end{document}
